\chapter{爲什麼要與人妻交往}

\section{與人妻交往的經濟學分析}

很多人把女人比作汽車。有五菱宏光，有坦克300，有豐田花冠，有賓士S級。對性取向正常的男性來說，對女人有需求是正常的，但是擁有女人的方式也多種多樣。就像如果你必須要用汽車，你可以買新車、可以買二手車、可以租車、可以打車、可以拼車。而與人妻交往，就相當於你直接去隔壁老王家把他家的汽車開走了。

現在中國和所有西方國家的男女關係，用車來比喻，就是一個極其混亂而低效的汽車市場。首先是新車銷售價越來越高。彩禮幾十萬，買房加名的話，可以到幾百萬上千萬，酒席也要幾萬幾十萬。然後在這之前，還有長期的考驗，考驗過程都要花錢。

你以爲你花了這些錢能買到一輛新車，結果人家給了你一個爆缸、爆胎過的二手車、事故車，做了板金，然後還按照新車的價格賣給你。如果你發現用新車價買了二手車了想要找人理論，他們還會倒打一耙，給你扣個「新車情結」的帽子，說什麼二手車開起來和新車一樣，還說我歧視二手車、事故車，真正優秀的賽車手不會嫌棄二手車等等讓人聽不懂的話。

更可怕的是，你花錢買完車，開了兩個月，某天回家車沒了。結果發現你的汽車自己啓動了自動駕駛模式開走了。然後你的手機收到一張離婚協議書，叫你籤名。這時候你只能去小紅書發一篇筆記：「有沒有靠譜的北京/上海/深圳離婚律師推薦？」，正文是「結婚兩個月老婆拿走78萬彩禮，找某知乎李姓大V在境外購入一個比特幣轉走了，然後要求分走北京星河灣/上海檀宮/深圳波託菲諾純水岸6000萬豪宅折現3000萬現金。」

然後你想着，這種情況下，只有傻子才買車。所以你有出行需求就想着打車、租車。結果一打車，你就被拘留了。拘留的時候又說了一堆你聽不懂的話，什麼你打車是侵犯汽車權益、侵犯汽車尊嚴、物化汽車了。還說什麼有一些車出售的時候說是自駕車，結果被司機強迫出來開網約車了。等等等等。

在這種情況下，其實從經濟上說，對大部分人來說，最有性價比的，就是睡別人老婆。購車款彩禮別人出，房貸車貸別人出，車險醫療保險別人出，你只要負責去享受就行了。不小心搞大人妻的肚子，還有「鐵綠帽法」\footnote{指在中國、法國等國家，丈夫有責任供養妻子婚內出軌所生子女的法律。讓你戴的這個綠帽像鐵帽子一樣，永遠摘不下來。}兜底，不用負責。而且一切完全合法，也符合現代價值觀（見後文），何樂而不爲呢？

而且，我們不僅要找人妻，還要找有錢人的老婆。有錢人的老婆一般都是我們一般無法負擔得起的昂貴美女。但是當她們成爲人妻之後，按照這本指南操作，可能不花錢直接享用在這個時代本來專屬有錢人的美女。還有什麼比睡有錢人老婆更快樂的事情呢？

\section{與人妻交往的法理分析}

首先，雖然我不是律師，但是我的法律能力絕對超過絕大多數律師。我曾經起訴美國總統，在美國聯邦法院和美國聯邦檢察官辯論不佔下風。而且更重要的是，我在這裏說的絕對是真話實話。

記住，不要相信抖音等任何中國的平臺的所有的「律師普法影片」。在中國，只有老老實實地讀法條和判例才有意義。（雖然中國不是判例法國家，但是從判例的分布可以倒推法院對法條解讀方式的機率分布。）中國大陸平臺對普法影片的審核標準只有一個：維穩。就是只要能達到維穩的目的，律師胡言亂語的影片也是可以發出來的。而如果一個律師，僅僅講述一個真實的判例，只要不利於「維穩」，也可能立刻被刪。比如中國有很多女的結婚後不讓碰，然後幾個月甚至幾天後就跑路離婚，然後彩禮一分不退的案例。但是你如果要在抖音講這些案例，是發不出來的。所以如果你靠抖音而不是法條或者科學取樣出來的案例，你是無法了解到真實的法律的。

\subsection{出軌的合法性}

\textbf{在中國和所有西方國家，人妻婚內出軌都是合法的。}出軌合法的核心是，女人有着完全的性自主權和生育自主權。無論你是女人的父母、男朋友、配偶、情人、子女、阿訇、校長、老闆、黨支部書記，你都沒有權利強制干預任何一個女人的身體自由和生育自由。所以包辦婚姻、丈夫強制要求妻子性愛、生育之類的行爲，是100\%違法的行爲。甚至可以判定爲屬於強奸罪、強奸未遂之類的嚴重暴力犯罪。

郭菊陽說服我跟她上牀的時候，說過她結婚是被她爸脅迫的，她內心最愛的其實是我。那麼如果她沒有做僞證的話，按照中華人民共和國法律，她老公應該是犯了強奸罪，應當處以3年以上10年以下有期徒刑。然後她爸屬於強奸罪的共犯，應當處以1年以上3年以下有期徒刑。如果郭菊陽是做僞證，那麼應當處三年以下有期徒刑或者拘役。

同樣的，\textbf{在中國和所有西方國家，黃毛和人妻約會都是合法的。}首先，一個婦女，是有完全的身體自主權的。也就是說，一個婦女的陰道和子宮，都不是任何人的所有物。不是父親的，也不是配偶的，更不是網友的。所以黃毛和人妻約會，不侵犯任何人的任何法律保護的權益。反而黃毛給人妻提供了在婚內無法體驗的快樂體驗，給別人帶來了額外的收益。

也就是，任何人都沒有對人妻身體的佔有權、支配權，但是人妻有着100\%追求自己幸福的權利，所以黃毛和人妻約會也是完全合法的。

前面說的，黃毛和人妻交往、約會是完全合法的。那麼自然，苦主也是沒有任何權利對人妻或者黃毛進行任何暴力行爲。

黃毛與人妻的生命權和健康權，比人妻的性自主權還要高，更不用說苦主對人妻的幻想中不存在的控制權。苦主對黃毛和人妻並不存在對其他人不存在的施暴的權利。

雖然中國婚姻法規定，夫妻之間有忠實義務，但是這種義務僅僅是原則性的，即使違反了，也只是在離婚時多判一點財產。而這的優先級其實是屬於最後的。最高人民法院的法律解釋說明，任何人無法僅僅依靠這一條義務進行維權。

而各種判例也是一直遵循這種原則。各種離婚判例，如果苦主對出軌的人妻施暴，在離婚中，家庭暴力的重要程度遠遠大於所謂違反忠實義務，財產分割都會偏向於女方。

\subsection{性自由的歷史根基}

前面說到，黃毛和人妻約會的行爲，無論是對黃毛還是人妻，無論是在中國，美國還是其他任何西方國家，都是完全合法的。這個原因就是，身體和性愛自主權，是根植於法國大革命和美國獨立戰爭，歷經美國南北戰爭、巴黎公社、辛亥革命、俄國十月革命、二戰反法西斯戰爭、中國人民的偉大的解放戰爭、還有美國民權運動，全球一波又一波的革命之後，確立下來的基本原則。

法國大革命的人權宣言寫道「人生而自由」。1950年新中國第一部法律《婚姻法》發表，毛主席盛贊爲「普遍性僅次於憲法的根本大法」。

也就是說，這種女性解放和自由，還有和人妻約會的合法性，是憲法水平、甚至是全世界所有文明國家的立國根基水平的。法律的保護是來源於最底層的。除非再經歷一場二戰水平烈度的戰爭，不可能有所改變。

\subsection{婚姻的法律起源}

首先，中國古代的婚姻制度，在現在中國可以說是從文化上、法律上、和道德上都完全斷絕傳承了，除了共用「婚姻」這個漢語詞，和現在的婚姻沒有關係。

現在無論是東亞還是西方的婚姻制度都是源自中世紀天主教的婚姻制度。中世紀天主教是一夫一妻制。天主教的婚姻是神聖的，婚姻必須基於愛情。這種愛情不只是現在的見色起意發情，而是一種神聖的感情。也就是夫妻雙方對對方，都需要像愛上帝一樣全身心地愛對方，夫妻雙方都不能和配偶以外其他人有性行爲，褻瀆婚姻相當於瀆神。然後在中世紀天主教婚姻制度下，離婚基本上都是不被允許的。天主教中，這種夫妻之間的愛情，是高於繁衍和後代的。甚至婚姻本身和繁衍後代無關，所以即使無子嗣也不能離婚。

到16世紀，英國國王亨利八世（就是那個娶了6個老婆砍了兩個老婆的殺妻狂魔亨利八世）爲了能夠隨意離婚換老婆，宣布脫離天主教會，創立英國國教會，讓自己成爲英國國教會的元首。自此英國的婚姻制度，就成爲了一個可以隨意離婚的天主教式婚姻。

而這種新教式婚姻制度，也就是這種臨時補丁式的可離婚式天主教婚姻，隨着西方文化的全球擴張，慢慢佔領了全世界，包括中國。這種過渡性的婚姻制度，其實很可笑，就是你可以隨便多次更換配偶，但是在兩次更換的中間，要像愛上帝一樣全身心地愛配偶——這本身就是違反理智和博弈論的一種可笑的制度。

而之前說的性自主權，出軌合法的觀念，又隨着法國大革命、十月革命、中國革命，作爲一個高於婚姻的人權概念，被各國所接受。這樣又進入了一種比較合理的制度，也就是無論是否結婚，性行爲和生育都不能違反女性的自由。如此一來，自由主義的邏輯又完全自洽了。

\subsection{離婚和鐵綠帽法}

但是性自主權，又是完全獨立於婚姻之上的。所以各國的離婚法律，法律精神的基礎，還是基於那一套可離婚的天主教一夫一妻制。而中國的婚姻和離婚制度，在天主教的基礎上，又加上了一套計劃生育的道德觀。在計劃生育的價值觀內，生育越多越邪惡，生育越少越光榮。所以在現代中國，那種沒有生育的愛情，婚姻是最偉大的。也就是說，最偉大、最政治正確的中國愛情故事，就是李松堅和凌菲菲的愛情故事了。一個億萬富豪愛上一個二婚帶男孩的女人，然後爲了愛情結合，億萬富豪把一套5億的檀宮別墅送給了這個女人。之後這個二婚帶男孩的女人又因爲愛情，愛上了一個保安隊長，然後勇敢地追求愛情，和富豪離婚，和保安隊長私奔追求愛情去了。

理解現代中國和西方婚姻的天主教背景，也就很好理解中國和西方的離婚制度了。也就是說只要一結婚，妻子都是丈夫的神。不管多久離婚，不管有沒有小孩，不管妻子有沒有出軌，只要丈夫比妻子富有，只要丈夫收入比妻子高，都需要大出血。

在美國，有一種制度叫做alimony，可以翻譯爲離婚贍養費。類似中國的離婚撫養費。Alimony的出發點是，如果沒離婚，那麼丈夫的收入都是妻子的，爲了保證公平，離婚後，高收入的丈夫必須向低收入的妻子支付收入的一部分，保證妻子不會因爲丈夫主動離婚而降低生活水平。當然美國的法律制度還是很注重男女平等的，如果妻子收入高於丈夫，也是一樣。

而生育，子女，血緣在天主教、基督新教的婚姻制度裏面是無關緊要的。所以子女撫養費，大部分時候都低於離婚財產分割甚至alimony。而妻子婚內出軌生下的小孩，默認也是丈夫的兒子，丈夫也是有必要一直支付撫養費的。

在美國的法律體系下，雖然綠帽是默認，但是丈夫如果不想認婚內非生物學子女的話，可以提起親子鑑定並聲明脫離父子關係。但是在法國，男性無權脫離婚內妻子出軌所產下子女的父子關係，必須一直支付撫養費，這就是我稱爲的「鐵綠帽法」——像滿清的鐵帽子王一樣，永遠也摘不下。

而中國也是實行鐵綠帽法，只不過在鐵綠帽法之上實行和稀泥法。婚內妻子出軌生下的非親生子女，或者繼子、繼女，在離婚後，男性也是有撫養義務的。但是中國特色和稀泥下，如果男方非常堅定嚴正抗議，或者輿論反撲比較大，法院會通過離婚財產分割偏向等方式補償男方。也就是說，任何一個進入婚姻的男人，最好掂量一下自己在法院的能量和對戴綠帽幫別人養小孩的接受程度，因爲中國默認是實行鐵綠帽法的。

\section{與人妻交往的道德分析}

其實通過前面的各種分析，其實與人妻交往的道德性也是顯然的了。

首先，很大一部分單身男性，因爲成本太高無法結婚甚至無法有性生活。他們有權利用合法的方式滿足自己的生理需求。

另外，很多女人在婚姻生活中，其實是很痛苦的，她們因爲害怕一些過時的、法律不支持的、不合時宜的傳統，頂着自相矛盾的道德，被迫和一個不喜歡的男人在一起生活，而如果人妻能和黃毛交往，她們也能夠很快樂。

而唯一似乎受損的苦主們，他們以爲受損的權利，真的有任何支持嗎？有任何法律還是道德，認爲一個女的只要領證，就要成爲配偶的附庸，喪失自由嗎？

那麼，我們和人妻交往約會，又有什麼道德缺陷呢？

如果我們找的是有錢人的老婆，那就更沒有道德缺陷了。在這個所謂自由市場經濟的時代，權力和資本勾結，有錢人、老人壟斷了所有的社會資源。年輕男性讀書、工作、創造了科技、網際網路、AI等所有財富，卻被中老年男人和富二代們奪去了絕大部分的收益，踩在我們的身體上，佔有了所有的女人。然後還反過頭來跟我們說要繼續奮鬥。

美國皮尤研究中心（Pew Research Center）研究表明，現在美國35歲以下年輕人佔有社會財富總量的比例，一直在刷新美國200多年歷史的最低點。而中國、日本、歐洲也無不如此。在這種幾百年來對年輕人最不公平的社會之中，通過偉大的愛情，對有錢人的老婆進行再分配，是再正義不過的事情了。

\section{中國現代的價值觀}

上面的法理和道德分析，用的是中國和西方法律上實質的現代自由主義。自由主義即使實際落地可能會造成一些問題，但是至少是邏輯自洽的。中學生可以隨便戀愛上牀，人妻也能隨便出軌。而儒家思想，其實邏輯也是自洽的。中學生15虛歲的七夕節出花園\footnote{潮汕少男少女的成人禮。}過後就可以結婚了。

但是在真實的自由主義法律體系下，現在中國臺面上的所謂主流的政治正確的價值觀，我稱之爲「保進步守」。也就是極端保守地擁護大約100年前的進步主義。這套進步主義，可以說是一套中華農民劣化版的普魯士-納粹總體戰價值觀。婦女解放反而是附帶。也就是說，無論男女，必須放棄一切的感情、娛樂，甚至要放棄家庭生活和生育，而把一切都投入到「工業發展」和「中華民族的偉大復興」。所以需要封禁索尼PS2遊戲機，需要禁止早戀，需要禁止一切和高考無關的東西，需要計劃生育。所有現實生活中無窮的苦，都是靠承諾的「幸福的彼岸」來支撐。你想想你高考前是不是有着無數的彼岸、天堂的承諾和地獄的恐嚇，和世間所有邪教又有何區別？理論上，這種普魯士-納粹式的進步主義敘事需要通過一場總體戰，獲得超額的收益，否則一整代人就白白犧牲了——最後文化停滯、全民創傷集體爆發、人口雪崩民族消亡。

而在六四事件之後，隨着所有的理想和信念消亡，普魯士價值觀繼續墮落到全世界最低劣的拜金主義、拜物主義價值觀。什麼精神追求都是政治錯誤的，那麼那些沒有堅持的人，也就只剩崇拜金錢了。再結合普魯士價值觀，雙重惡心：禁止早戀、禁止娛樂、計劃生育不是爲了偉大復興，也不是爲了對美/日/俄的復仇總體戰，而是爲了金錢。這種價值觀也迅速讓上上下下全體腐化。到現在爲止，中國女人人均郭菊陽，也就是可以犧牲一切，十分努力地用身體換取金錢的「奮鬥雞」。而中國男人，則是人均相信「贏了會所嫩模」的「奮鬥嫖客」了。

我有一個在北京微軟工作的朋友，曾經在北京多次重複嫖宿一個湖南省出身的18歲妓女。當時正值房產泡沫時期，每一次需要5000人民幣。有一天事後，他抱着湖南雞，像所有的男人一樣，開始勸雞從良。他忍不住詢問，她長得這麼漂亮，怎麼淪落到北京做雞呢？她很認真地回答，她這麼努力是爲了在北京買房，爲了有一個自己的家。她開始算起來：每一次接客，被中介抽水後可以賺3000多塊，然後趁年輕漂亮一天努力多接幾單，幹到二十六七歲就能夠在北京買到一套1000萬左右的房子了。說完，我朋友爲她努力做雞買房的愛國行爲感動萬分\footnote{我這個朋友，因爲看過我在2011-2012年寫的關於土地財政和地方債的文章，所以就堅決不買房而一直租房。實在是「不愛國」。}，決定再操她一次。經過這麼一輪交心的交談，他們兩人，又幹得更用力了……我的朋友沉浸在18歲湖南雞誘人的肉體之中，而湖南雞離在北京買房的夢想又近了一步。這就是美好的愛情吧。

我在2021年第一次在深圳把郭菊陽約出來的時候，她和我一見面，就開始炫耀她在波託菲諾純水岸15期4棟的那套價值6000多萬的豪宅了。當時我看着她那既熟悉而又陌生的臉，像在《了不起的蓋茨比》裏面一樣，感覺her voice is full of money（她的聲音充滿了金錢）。她看我毫無反應，又強調了一遍她的房子價值6000多萬。當時她那驕傲的表情，像是一個已經成功到達目標的湖南雞的MVP結算畫面，不對，像是一個至少3倍超額完成目標的湖南雞的MVP結算畫面。說着說着，她又開始罵附近白石洲城中村還是哪裏的什麼建設計劃，會影響她的房產價格。因爲我實在是不感興趣，左耳進右耳出，也忘了具體她當時罵的是什麼了。或許也是因爲她雞成這樣，才讓我和朋友一樣想要勸雞從良吧。

但是話說回來，價值觀是價值觀，法律是法律。真的有人按着頭逼着你奮鬥，逼着你拜金嗎？在學校的時候泡泡妞，上班的時候摸摸魚，有錢就到處旅遊喫喫喝喝，上班時間泡泡家庭主婦，下課的時候就泡泡中學生，紐約香港住膩了，就我一樣買張機票在東京買套房旅居\footnote{那些在北上廣深買房的，隨便賣一套就可以在大阪和東京（港區以外）住得很舒服了。}，寫寫書，追求一下藝術理想。這樣的生活，難道不是開心多了？
